% PROPOSED response to Reviewer1, Comment1 ONLY. Original manuscript untouched.
\noindent\textbf{Reviewer1, Comment1.}
The sample-adaptive construction does not establish whether equally or nearly
suitable references can materially change DI for the same observed sample.

\medskip
\noindent\textbf{Proposed response.}
We agree that sample-adaptive initialization is not a proof of invariance to
alternative references. Consistent permutations of cluster labels leave all
translations and DI unchanged, but different partitions or medoid geometries
need not do so. Similar objective values alone do not establish a DI bound.

We performed a focused check on the existing clustered population with $N=144$,
$n=8$, unequal IPs, and two fixed observed samples chosen before inspecting
sensitivity. For each sample we used ten runs of the official sample-initialized
procedure and ten separately labeled alternative-initialization stress tests.
Population, samples, IPs, uniform box kernel, covariance-based Scott bandwidth,
and angular scaling were fixed. The full OT, refinement, hardening, squared
medoid and matching pipeline was recomputed; no DI values or densities were
fabricated. A separate audited adapter executes the package routines while
respecting the manuscript-specific choices where the public implementation
differs. The original DI formula was retained.

The manuscript states a preliminary entropic OT objective but no numerical
equivalence tolerance or final-reference minimization criterion. We therefore
report both its native normalized loss and refined soft compactness, actual
hard-cluster IP totals, and comparisons at 0.1\%,1\%, and 5\% loss tolerances.
The official runs produced two distinct references per sample with identical
preliminary losses; the DI ranges were 0.0279 and 0.0043. In the stress test, an
alternative with lower native OT loss and lower refined compactness loss had
gaps 4.45\% and 2.65\%, respectively, but changed the second sample's DI from
$-0.58223930146083458$ to$-0.42201965129016566$. Both compared TSP paths were
optimal for their preliminary centers, and $D_{\mathrm{net}}$ remained negative
and away from zero. This substantial contrast therefore does not arise from a
near-zero sign reversal or merely a poor TSP path.

The tolerance matters: the stricter joint 0.1\% and 1\% rules admitted no distinct
pairs in this set. Moreover, because the manuscript recomputes the cost median,
using a common median for objective comparison makes the alternative's OT loss
20.86\% lower, rather than nearly equal. We do not interpret the example as a
universal statement about arbitrary near-optimal references. It shows why fixing
the construction provides a reproducible convention, not a general robustness
guarantee. The proposed manuscript clarification therefore qualifies DI as
reference-conditional, explicitly distinguishes official randomness from the
initialization stress test, and restricts the numerical conclusion to the
examined population and samples. The compact table and one figure provide the
computed references and density evidence; full settings, seeds, coordinates,
and verification are supplied separately.
